\chapter{基于空间信号表征的多目标协同鲁棒定位方法}
\label{ch:multi_target_localization}

\section{引文}
\label{sec:4_1}

本章将复杂城市环境下的多目标并发定位问题重构为空间信号场表征与区域级判别的联合建模问题，并在此基础上进一步考虑多无人机协同条件下的节点掉点与残缺观测问题。与传统方法主要依赖离散峰值检测、逐目标分离或局部判定不同，本章从问题建模和能力迁移两个层面建立统一方法框架，以同时应对稀疏扫描、多径鬼影和协同网络不完整带来的定位退化。

\begin{figure}[htbp]
    \centering
    \includegraphics[width=0.94\textwidth]{ch4/第四章图.png}
    \caption{第四章多目标协同定位方法总体思路}
    \label{fig:ch4_overall_innovation_framework}
\end{figure}

如图~\ref{fig:ch4_overall_innovation_framework} 所示，本章的研究内容主要包括两个方面。第一，面向完整或相对完整观测条件下的多目标并发定位问题，针对无人机空间扫描带来的离散采样、局部多峰和多径鬼影影响，引入 Radio Map 作为中间空间表征，将离散峰值检测问题转化为空间信号场成像与区域级判别问题，以实现多目标系统的高保真建模。第二，面向多无人机协同过程中的节点掉点场景，构建完整观测教师网络与残缺观测学生网络之间的知识迁移框架，将完整协同观测条件下获得的空间判别能力迁移到残缺观测场景，从而提升掉点条件下的多目标协同定位鲁棒性。基于上述两方面内容，本章前半部分首先围绕多目标并发定位的场景建模、稀疏采样特性与信号成像展开分析，后半部分再进一步讨论节点掉点条件下的知识蒸馏与鲁棒协同定位问题。

\section{基于无人机空间扫描的多目标定位问题描述与信号成像建模}
\label{sec:4_2}

本节面向完整或相对完整观测条件下的多目标定位问题，依次讨论任务定义、无人机空间离散采样、噪声扰动、局部多峰以及多径伪峰对传统峰值判定的影响，并在此基础上引出后续 Radio Map 信号成像与区域级判别思路。

\subsection{多目标并发定位场景与任务定义}
\label{subsec:4_1_1}

本章研究的对象为复杂城市环境下地面非法低轨卫星终端的多目标并发定位问题。与第三章中面向单目标连续跟踪的定位任务不同，本章关注的是：在同一观测时刻或同一空间区域内，存在多个非法终端同时活动时，如何依托低成本无人机平台完成对多个被定位目标的联合识别与位置判定。这里，无人机平台并不是被定位对象，而是承担空间扫描、信号采集和定位解算任务的机动观测平台；地面非法终端则是定位系统需要识别和估计位置的辐射源目标。

图~\ref{fig:ch4_multi_target_uav_scenario} 给出了本章对应的多目标并发定位场景示意图。

\begin{figure}[htbp]
    \centering
    \includegraphics[width=0.84\textwidth]{ch4/1_多目标并发定位场景示意图.jpg}
    \caption{基于低成本无人机平台的地面多目标非法终端并发定位场景}
    \label{fig:ch4_multi_target_uav_scenario}
\end{figure}

从任务形式上看，多目标并发定位与单目标定位存在本质差异。单目标场景下，系统主要关注单个目标位置的连续估计，其核心问题通常表现为轨迹误差抑制与状态更新稳定性；而在多目标场景下，系统不仅需要判断目标是否存在，还需要同时回答有多少个目标和每个目标位于何处两个问题。设在给定观测区域内共有 \(K\) 个同时活动的地面非法终端，其真实位置集合记为
\begin{equation}
\mathcal{P}
=
\left\{
\mathbf{p}_1,\mathbf{p}_2,\ldots,\mathbf{p}_K
\right\},
\qquad
\mathbf{p}_k=[x_k,y_k]^T,\; k=1,2,\ldots,K,
\label{eq:4_1_target_position_set}
\end{equation}
其中，\(K\) 表示当前区域内真实目标数量，\(\mathbf{p}_k\) 为第 \(k\) 个目标的二维平面坐标。面向多目标场景的区域级定位判别任务，是指无人机平台根据在空间扫描过程中获得的无线观测数据，对目标区域形成判别结果并给出对应的位置估计。本文当前实现重点在于多目标场景下的区域级判别与位置推断，而非完整目标数估计与实例级检测。

在本文场景中，无人机平台通过飞行观测方式对
目标区域执行空间扫描。受飞行航迹、驻留时间、传感器视场和载荷能力等因素限制，无人机不可能在整个区域内连续、同步地获得全域信号场，而是只能在有限位置获取离散观测样本。也就是说，本章所处理的多目标定位问题，从观测层面上首先表现为基于无人机机动采样的空间离散感知问题。与固定基站阵列或密集地面测量不同，这种观测方式具有较强的灵活性和区域覆盖能力，但同时也引入了采样稀疏、观测不均匀和局部信息缺失等问题。

进一步地，复杂城市环境中的高层建筑、街道走廊、玻璃幕墙以及其他反射界面，会使无人机接收到的信号同时包含真实目标响应、多径反射分量和背景噪声扰动。对于单目标场景，这些因素主要表现为位置估计偏差；而在多目标场景中，上述因素还会对区域判别结果产生直接影响。当多个地面终端同时存在时，不同目标的空间响应可能发生局部重叠，噪声扰动可能造成局部伪峰，多径反射还可能在非真实位置形成附加热点。由此可见，面向多目标场景的区域级定位判别并不是若干单目标定位结果的简单叠加，而是一个同时涉及区域级识别、目标区分和位置推断的联合问题。

% 从工程实现角度看，传统多目标无线定位方法通常依赖逐目标搜索、逐目标分离或迭代残差消除等串行处理方式。该类方法在目标数量较少时仍可工作，但当目标并发数增加时，其计算负担通常随目标数上升而增加，同时在强多径和高噪声条件下容易出现漏检与虚警并存的情况。对于低成本无人机平台而言，这种处理方式难以同时满足复杂环境下的实时性和稳定性要求。因此，本章的任务并不只是提高单个目标的位置估计精度，更重要的是建立一种适合无人机机动观测场景的多目标联合判定方法，使系统能够在有限计算开销下，对目标区域内多个非法终端的空间分布进行有效识别。

基于上述分析，本章关注的问题可进一步表述为：无人机平台在目标区域上空或周边执行空间扫描任务，通过采集由多个地面非法终端共同形成的离散无线观测，构建目标区域的空间信号分布表征，并在此基础上完成多目标的并发识别与定位。下文将进一步说明，若仍直接在离散观测点或局部峰值层面进行判定，容易受到稀疏采样、噪声和多径干扰影响；因此，有必要先将无人机采集到的离散观测转换为具有区域结构的信号图像，再在图像域内开展区域级判别。

\subsection{无人机稀疏采样、噪声扰动与局部多峰问题分析}
\label{subsec:4_1_3}

在基于无人机平台的多目标并发定位任务中，目标区域的空间信号分布只能由无人机沿飞行航迹在有限位置上离散采样获得。受飞行高度、航迹长度、悬停时间、采样周期以及平台载荷能力等因素约束，无人机通常难以完成高密度、全覆盖的连续测量，因此本章首先面对的是空间稀疏采样条件下的多目标定位问题。

设无人机在目标区域内共获取 \(N\) 个采样点，其空间坐标与对应观测值记为
\begin{equation}
\mathcal{D}
=
\left\{
(\mathbf{s}_i,r_i)
\right\}_{i=1}^{N},
\qquad
\mathbf{s}_i=[x_i,y_i]^T,
\label{eq:4_2_sparse_sampling_data}
\end{equation}
其中，\(\mathbf{s}_i\) 表示第 \(i\) 个采样点的位置，\(r_i\) 表示对应位置处的接收信号强度观测值。对于包含 \(K\) 个目标的场景，\(r_i\) 一般可写为
\begin{equation}
r_i
=
\sum_{k=1}^{K} h_k(\mathbf{s}_i;\mathbf{p}_k)
+
n_i,
\label{eq:4_3_multi_source_observation}
\end{equation}
其中，\(h_k(\mathbf{s}_i;\mathbf{p}_k)\) 表示第 \(k\) 个真实目标在位置 \(\mathbf{s}_i\) 处产生的空间响应，\(\mathbf{p}_k\) 为第 \(k\) 个目标的真实位置，\(n_i\) 为观测噪声项。由式~\eqref{eq:4_3_multi_source_observation} 可知，每一个离散采样值都是多目标响应与噪声的叠加结果，而非某一单独目标的纯净观测。

若进一步考虑理想路径损耗模型，则单个目标的空间响应可近似表示为
\begin{equation}
h_k(\mathbf{s}_i;\mathbf{p}_k)
=
A_k
-
10\eta \log_{10}\!\left(d_{ik}\right),
\qquad
d_{ik}=\|\mathbf{s}_i-\mathbf{p}_k\|_2,
\label{eq:4_4_path_loss_model}
\end{equation}
其中，\(A_k\) 表示第 \(k\) 个目标的等效发射强度，\(\eta\) 为路径损耗指数，\(d_{ik}\) 为采样点与目标之间的欧氏距离。式~\eqref{eq:4_4_path_loss_model} 仅给出理想条件下的距离衰减趋势，复杂城市环境中的遮挡、散射和反射会使实际观测偏离这一关系。

若仅存在单一真实目标且噪声较弱，则连续空间内的响应通常表现为围绕目标中心逐渐衰减的单峰结构。设理想连续响应记为 \(g(\mathbf{x})\)，噪声扰动记为 \(\xi(\mathbf{x})\)，则实际观测场可表示为
\begin{equation}
\tilde{g}(\mathbf{x})
=
g(\mathbf{x})+\xi(\mathbf{x}).
\label{eq:4_5_noisy_field}
\end{equation}
在理想条件下，\(g(\mathbf{x})\) 往往仅在目标附近存在单一主极值点；但当 \(\xi(\mathbf{x})\) 在局部区域内产生较强波动时，\(\tilde{g}(\mathbf{x})\) 可能在单一目标附近分裂出多个相邻极值点。对离散采样情形，这一现象可进一步表示为
\begin{equation}
\hat{r}_i
=
r_i+\delta_i,
\label{eq:4_6_sample_perturbation}
\end{equation}
其中，\(\delta_i\) 表示由测量抖动、采样误差和随机噪声引入的附加扰动。当部分采样点的扰动明显偏大时，原本由单一目标形成的局部能量团簇便可能在离散观测图中表现为多个分裂高点。

为了说明这一点，可将局部极值集合记为
\begin{equation}
\mathcal{M}
=
\left\{
\mathbf{m}_1,\mathbf{m}_2,\ldots,\mathbf{m}_L
\right\},
\qquad
\mathbf{m}_l \in \Omega,
\label{eq:4_7_local_maxima_set}
\end{equation}
其中，\(\Omega\) 为观测区域，\(\mathbf{m}_l\) 表示第 \(l\) 个局部极大值位置，\(L\) 为局部极值点总数。在理想单目标条件下，通常有 \(L \approx 1\)；而在稀疏采样和噪声扰动作用下，可能出现 \(L>1\) 的情况，即同一真实目标对应多个局部峰值。

对于无人机机动采样平台而言，上述问题还会因采样分布不均匀而进一步加剧。不同区域的采样密度与相对观测角度并不完全相同，局部样本缺失会削弱对真实空间响应边界的刻画能力，使观测场更容易受异常样本支配。在多目标并发场景中，这意味着多个局部峰值未必对应多个不同终端，它们也可能来自同一目标在噪声和采样扰动下形成的分裂响应。若系统直接依据峰值点数完成目标判定，便会将同一真实目标误识别为多个候选点。因此，问题的关键不只是寻找空间高值点，而是判断若干局部高值是否属于同一连续响应区域，判定单元需要从单点极值转向区域结构表达。图~\ref{fig:ch4_sparse_sampling_multipeak} 给出了无人机稀疏采样与噪声扰动条件下局部多峰现象的示意图。

\begin{figure}[htbp]
    \centering
    \includegraphics[width=0.84\textwidth]{ch4/图片.png}
    \caption{无人机稀疏采样与噪声扰动条件下的局部多峰现象}
    \label{fig:ch4_sparse_sampling_multipeak}
\end{figure}

\subsection{多径伪峰干扰下传统峰值判定方法的局限性分析}
\label{subsec:4_1_4}

前一小节已经说明，稀疏采样与噪声扰动会使局部峰值个数偏离真实目标数目；进一步地，复杂城市环境中的多径传播还会引入鬼影干扰。对于无人机平台而言，建筑外立面、玻璃幕墙、街道峡谷和金属设施等反射体普遍存在，目标信号在传播过程中往往会经历反射、绕射和散射，因此无人机接收到的是多条传播路径叠加后的复合响应。这样一来，峰值集合中会同时混杂真实目标主峰、同源分裂峰以及多径鬼影峰，若仅依据峰值幅度或局部邻域极大关系进行判定，便难以准确区分这些结构。

为了说明这一局限，可设传统峰值判定规则为
\begin{equation}
\mathbf{x}\in\mathcal{T}
\quad \Longleftrightarrow \quad
r(\mathbf{x}) > \tau
\ \text{且}\
r(\mathbf{x})=\max_{\mathbf{y}\in\mathcal{N}(\mathbf{x})} r(\mathbf{y}),
\label{eq:4_12_peak_detection_rule}
\end{equation}
其中，\(\tau\) 为幅值阈值，\(\mathcal{N}(\mathbf{x})\) 表示 \(\mathbf{x}\) 的局部邻域，\(\mathcal{T}\) 为被判定为目标候选点的集合。该规则本质上只利用了是否超过阈值和是否为局部极大值两类信息，而没有刻画高值点所属区域的整体纹理、扩散结构和空间连续性。因此，只要鬼影峰或噪声峰满足式~\eqref{eq:4_12_peak_detection_rule}，便会被纳入候选目标集合，导致虚警产生。

对于无人机平台的多目标扫描任务，这一问题还具有两个特点。其一，鬼影位置和强度与无人机观测视角、建筑布局以及目标分布共同相关，因此不同飞行位置下鬼影的空间表现并不稳定；其二，当多个真实目标同时存在时，不同目标的主响应与伪峰响应可能发生局部重叠，使得仅依赖峰值强度排序的方法更难完成正确区分。更重要的是，真实目标与伪峰的差异主要体现在区域层面的形态特征上，而非单个峰值幅度。由此可见，传统峰值判定方法的根本局限不在于阈值设置，而在于判定单元过于局部，无法利用空间响应的整体形态信息。

基于上述分析，本文不再将多目标并发定位问题视为离散观测点上的峰值搜索问题，而是将其重构为面向区域信号分布的判别问题，即先将离散 RSSI 观测映射为具有连续结构的 Radio Map，再在图像域内完成目标响应区域的结构化判别与位置推断。下一节将进一步给出离散采样信号成像建模与 Radio Map 构建过程。

\section{面向多目标定位的 Radio Map 构建}
\label{sec:4_3}

本节围绕面向多目标定位的 Radio Map 构建过程展开。整体上，本文以 OSM 导出的真实街区地理信息数据为底层场景来源，结合 MATLAB Site Viewer、Antenna Toolbox 与 RF Propagation Toolbox 构建复杂城市场景与传播仿真工具链。其中，OSM 数据用于提供道路与建筑几何信息，Site Viewer 用于场景导入、可视化检查与空间关系校验，Antenna Toolbox 与 RF Propagation Toolbox 分别负责收发天线参数描述和射线追踪传播计算。基于该工具链，本文获得覆盖目标区域 27,546 个网格点的高分辨率传播数据，为后续 RSSI 栅格映射和 Radio Map 构建提供基础输入。

在此基础上，本节首先给出离散采样信号成像建模与 Radio Map 的统一定义，随后介绍复杂城市传播场建模和 RSSI 栅格映射方法。图~\ref{fig:ch4_osm_siteviewer_toolchain} 给出了基于 OSM 与 MATLAB 工具链的复杂城市场景建模流程。

\begin{figure}[htbp]
    \centering
    \includegraphics[width=0.84\textwidth]{ch4/7_OSM与MATLAB工具链建模流程图.png}
    \caption{基于 OSM 与 MATLAB 工具链的复杂城市场景建模流程}
    \label{fig:ch4_osm_siteviewer_toolchain}
\end{figure}

\subsection{离散采样信号成像建模与 Radio Map 定义}
\label{subsec:4_2_1}

为将上一节提出的区域级信号表征思路落到可计算的空间模型上，本节首先对离散采样信号成像过程进行数学建模，并给出 Radio Map 的统一定义。设无人机在观测区域 \(\Omega\) 内获得的离散观测集合为
\begin{equation}
\mathcal{D}
=
\left\{
(\mathbf{s}_i,r_i)
\right\}_{i=1}^{N},
\qquad
\mathbf{s}_i=[x_i,y_i]^T,
\label{eq:4_13_discrete_samples}
\end{equation}
其中，\(\mathbf{s}_i\) 和 \(r_i\) 的含义与前文相同。本文首先在目标区域上建立二维规则网格
\begin{equation}
\mathcal{G}
=
\left\{
\mathbf{g}_{u,v}
\mid
u=1,2,\ldots,U,\;
v=1,2,\ldots,V
\right\},
\label{eq:4_14_grid_definition}
\end{equation}
其中，\(\mathbf{g}_{u,v}\) 表示第 \((u,v)\) 个网格单元对应的空间位置，\(U\) 和 \(V\) 分别为网格在两个方向上的离散分辨率。在此基础上，通过空间映射、插值和平滑操作，将离散观测集合 \(\mathcal{D}\) 转换为定义在网格 \(\mathcal{G}\) 上的连续场近似
\begin{equation}
\hat{R}(u,v)
=
\mathcal{F}\!\left(\mathcal{D},\mathbf{g}_{u,v}\right),
\label{eq:4_15_field_reconstruction}
\end{equation}
其中，\(\mathcal{F}(\cdot)\) 表示由离散采样到连续场的映射过程，\(\hat{R}(u,v)\) 为第 \((u,v)\) 个网格位置处的信号强度值。于是，整个观测区域的无线空间响应可进一步表示为二维矩阵
\begin{equation}
\mathbf{R}
=
\left[
\hat{R}(u,v)
\right]_{U\times V},
\label{eq:4_16_radio_map_matrix}
\end{equation}
式~\eqref{eq:4_16_radio_map_matrix} 即构成了本章后续所使用的 Radio Map 表征。

与原始离散采样点集合相比，Radio Map 的意义不只在于插值补全，更在于将目标响应从离散点提升为具有空间邻域关系的区域结构。对于同一真实目标，噪声诱发的若干局部高值点在热力图中通常会落入同一连续响应区域；对于多径鬼影，其空间形态、边界特征和邻域衰减方式则往往与真实目标不同。因此，将离散无线观测映射为 Radio Map，本质上完成了从数值域峰值判定到图像域区域级判别的表征转换。

进一步地，若记传统峰值判定的目标候选集合为
\begin{equation}
\mathcal{T}_{\mathrm{peak}}
=
\left\{
\mathbf{x}
\in \Omega
\mid
r(\mathbf{x})>\tau,\;
r(\mathbf{x})=\max_{\mathbf{y}\in\mathcal{N}(\mathbf{x})}r(\mathbf{y})
\right\},
\label{eq:4_17_peak_based_candidates}
\end{equation}
则本文所采用的图像域判定思路可表述为：在 Radio Map \(\mathbf{R}\) 上寻找满足目标区域特征约束的响应子区域集合
\begin{equation}
\mathcal{T}_{\mathrm{region}}
=
\left\{
\mathcal{B}_1,\mathcal{B}_2,\ldots,\mathcal{B}_{\hat{K}}
\right\},
\label{eq:4_18_region_based_candidates}
\end{equation}
其中，\(\mathcal{B}_m\) 表示第 \(m\) 个判定得到的候选目标区域，\(\hat{K}\) 为区域级判别输出的候选区域数量。与式~\eqref{eq:4_17_peak_based_candidates} 相比，式~\eqref{eq:4_18_region_based_candidates} 的判定单元由点集扩展为区域集合，这意味着目标识别过程不再只关注单个位置处的局部极值，而是同时考虑区域范围、局部纹理和邻域结构等多种信息。

基于上述建模思路，本章将多目标并发定位问题进一步转化为 Radio Map 上的区域级判别问题，即先将无人机采集到的离散 RSSI 观测映射为二维无线电热力图，再在图像域内对目标响应区域进行结构化判别并输出对应位置估计。图~\ref{fig:ch4_cross_modal_pipeline} 给出了从无人机离散观测到 Radio Map 成像，再到区域级判别的整体流程。

\begin{figure}[htbp]
    \centering
    \includegraphics[width=0.86\textwidth]{ch4/4_离散观测到RadioMap成像流程图.jpg}
    \caption{从无人机离散观测到 Radio Map 成像及跨模态区域判别的整体流程}
    \label{fig:ch4_cross_modal_pipeline}
\end{figure}

\subsection{基于射线追踪的复杂城市传播场建模}
\label{subsec:4_2_2}

为构建能够反映复杂城市环境传播机理的 Radio Map，本文首先引入基于射线追踪的城市传播场建模方法，对目标区域内的电磁传播过程进行刻画。多目标场景中的信号分布不仅由目标位置决定，还同时受到建筑物遮挡、街道走廊效应、玻璃幕墙反射以及非视距传播等因素影响，因此需要先获得具有物理一致性的空间传播响应。

在上一节所述工具链基础上，本文完成目标区域三维场景重建，并据此建立传播场模型。设目标区域为
\begin{equation}
\Omega \subset \mathbb{R}^2,
\label{eq:4_19_region}
\end{equation}
其对应的三维建筑集合记为
\begin{equation}
\mathcal{B}
=
\left\{
B_1,B_2,\ldots,B_M
\right\},
\label{eq:4_20_building_set}
\end{equation}
其中，\(B_m\) 表示第 \(m\) 个建筑物在三维空间中的几何实体，\(M\) 为场景内建筑物总数。集合 \(\mathcal{B}\) 为后续射线传播分析提供了几何基础。

在传播建模阶段，本文利用射线追踪模块计算目标终端至无人机观测平台之间的传播路径。对于给定目标位置 \(\mathbf{p}_k\) 与无人机观测位置 \(\mathbf{q}_i\)，接收信号可视为多条有效传播路径的叠加，即
\begin{equation}
r_{ik}(t)
=
\sum_{\ell=1}^{L_{ik}}
\alpha_{ik}^{(\ell)}
s_k\!\left(t-\tau_{ik}^{(\ell)}\right)
e^{-j2\pi f_c \tau_{ik}^{(\ell)}}
+
n_{ik}(t),
\label{eq:4_21_ray_sum}
\end{equation}
其中，\(L_{ik}\) 表示从目标 \(k\) 到观测点 \(i\) 的有效传播路径数，\(\alpha_{ik}^{(\ell)}\) 和 \(\tau_{ik}^{(\ell)}\) 分别表示第 \(\ell\) 条路径的复增益和传播时延，\(f_c\) 为载波频率，\(s_k(t)\) 为目标发射信号，\(n_{ik}(t)\) 为噪声项。式~\eqref{eq:4_21_ray_sum} 表明，接收信号不仅包含主路径分量，还显式包含由建筑反射、散射与绕射产生的附加路径分量，这些分量共同决定了最终空间信号场的分布结构。

在路径级传播结果基础上，本文进一步计算各观测位置处的接收功率，并将其作为后续 Radio Map 重构的原始输入。若记第 \(i\) 个观测点处来自第 \(k\) 个目标的接收功率为 \(P_{r,ik}\)，则其可写为
\begin{equation}
P_{r,ik}
=
\sum_{\ell=1}^{L_{ik}}
\left|
\alpha_{ik}^{(\ell)}
\right|^2
P_{t,k},
\label{eq:4_22_receive_power}
\end{equation}
其中，\(P_{t,k}\) 为第 \(k\) 个目标的发射功率。对于多目标场景，若目标区域内存在 \(K\) 个同时活动的地面非法终端，则无人机在位置 \(\mathbf{q}_i\) 处的总接收功率可进一步表示为
\begin{equation}
P_r(\mathbf{q}_i)
=
\sum_{k=1}^{K} P_{r,ik}
+
\varepsilon_i,
\label{eq:4_23_multi_target_power}
\end{equation}
其中，\(\varepsilon_i\) 表示背景噪声与测量误差项。式~\eqref{eq:4_23_multi_target_power} 给出了后续热力图构建所需的空间功率观测基础。

与经验路径损耗公式或简单空间插值相比，基于射线追踪的传播建模能够显式引入建筑几何结构、非视距遮挡和反射路径等环境因素，使得到的空间功率分布具有更明确的物理依据，并为多目标场景中的鬼影分析提供路径级解释。后续本文不直接对式~\eqref{eq:4_23_multi_target_power} 中的离散功率观测进行逐点判定，而是进一步将其映射到二维规则网格中，构建无线电热力图。

\begin{figure}[htbp]
    \centering
    \begin{minipage}[b]{0.78\textwidth}
        \centering
        \includegraphics[width=\textwidth]{ch4/11_城市峡谷RadioMap三维图A.png}
    \end{minipage}

    \vspace{0.7em}

    \begin{minipage}[b]{0.78\textwidth}
        \centering
        \includegraphics[width=\textwidth]{ch4/12_城市峡谷RadioMap三维图B.png}
    \end{minipage}
    \caption{复杂城市传播场条件下 Radio Map 的三维可视化结果（一）}
    \label{fig:ch4_radiomap_3d_gallery_1}
\end{figure}

\begin{figure}[htbp]
    \centering
    \begin{minipage}[b]{0.78\textwidth}
        \centering
        \includegraphics[width=\textwidth]{ch4/13_反射增强RadioMap三维图C.png}
    \end{minipage}

    \vspace{0.7em}

    \begin{minipage}[b]{0.78\textwidth}
        \centering
        \includegraphics[width=\textwidth]{ch4/14_阴影衰落RadioMap三维图D.png}
    \end{minipage}
    \caption{复杂城市传播场条件下 Radio Map 的三维可视化结果（二）}
    \label{fig:ch4_radiomap_3d_gallery_2}
\end{figure}

由图~\ref{fig:ch4_radiomap_3d_gallery_1} 和图~\ref{fig:ch4_radiomap_3d_gallery_2} 可以看出，在多目标并发辐射、建筑遮挡、街道走廊效应以及多径反射共同作用下，空间接收功率分布呈现出显著的非均匀峰谷结构。不同场景下的 Radio Map 在峰值位置、能量扩散范围以及局部起伏纹理上均存在明显差异，说明复杂城市环境中的空间响应具有显著的区域结构和场景依赖性，可为后续 RSSI 栅格映射和区域级判别提供更具判别性的输入基础。

\subsection{基于空间扫描观测的 RSSI 栅格映射方法}
\label{subsec:4_2_3}

在完成复杂城市传播场建模之后，本文进一步将无人机平台在空间扫描过程中获取的离散 RSSI 观测映射到二维规则网格上，为后续 Radio Map 连续化表达提供统一输入。射线追踪建模给出了具有物理一致性的空间传播响应，而后续区域级判别所需要的则是尺寸统一、坐标明确且便于结构化分析的二维灰度图像，因此有必要在二者之间建立一层栅格化映射过程。

沿用前文记号，设无人机共获取 \(N\) 个观测样本，第 \(i\) 个观测点在二维平面内的坐标记为 \(\mathbf{q}_i=[x_i,y_i]^T\)，对应 RSSI 观测值记为 \(r_i\)。

为了将无人机沿航迹获取的无序离散观测组织成适合图像处理的规则表示，本文在目标区域上建立二维规则网格。结合当前多目标验证场景，本文以 \(500\,\mathrm{m}\times 500\,\mathrm{m}\) 的室外复杂街区作为成像区域，并将该区域按照固定空间分辨率划分为规则栅格。每个栅格单元对应一个确定的空间位置，所有观测样本均根据其平面坐标投影到对应网格内。

在映射过程中，若某一栅格内仅包含一个观测点，则直接将其 RSSI 值赋给该栅格；若同一栅格内包含多个观测样本，则采用均值融合方式生成该栅格的初始信号强度值。对于未被无人机直接扫描到的网格单元，在本阶段不进行强行估计，而是保留为空值或零值，并在后续高斯平滑步骤中统一完成连续化补全。经过这一处理后，原始散点观测被统一组织为二维矩阵
\begin{equation}
\mathbf{S}
=
\left[
S(u,v)
\right]_{H\times W},
\end{equation}
其中，\(H\) 和 \(W\) 分别表示栅格在两个方向上的离散尺度，\(S(u,v)\) 表示第 \((u,v)\) 个栅格单元上的初始 RSSI 值。

记落入第 \((u,v)\) 个栅格单元的观测样本索引集合为 \(\mathcal{I}_{u,v}\)，则初始栅格值可写为
\begin{equation}
S(u,v)=
\begin{cases}
\dfrac{1}{|\mathcal{I}_{u,v}|}\displaystyle\sum_{i\in\mathcal{I}_{u,v}} r_i, & |\mathcal{I}_{u,v}|>0,\\[8pt]
0, & |\mathcal{I}_{u,v}|=0,
\end{cases}
\label{eq:4_24_grid_assignment}
\end{equation}
其中，空栅格先置为 0，并在后续平滑过程中完成连续化补全。为抑制局部孤立尖峰并增强空间响应的区域连续性，本文对初始矩阵 \(\mathbf{S}\) 施加二维高斯平滑，定义高斯核为
\begin{equation}
G_{\sigma}(a,b)=
\frac{1}{2\pi\sigma^{2}}
\exp\!\left(
-\frac{a^{2}+b^{2}}{2\sigma^{2}}
\right),
\label{eq:4_25_gaussian_kernel}
\end{equation}
其中，\(\sigma\) 为平滑核宽度参数。于是，平滑后的连续场表示可写为
\begin{equation}
\tilde{S}(u,v)=
\sum_{m=1}^{H}\sum_{n=1}^{W}
G_{\sigma}(u-m,v-n)\,S(m,n),
\label{eq:4_26_gaussian_smoothing}
\end{equation}
在此基础上，再对平滑结果进行线性归一化，得到最终用于后续分析的 Radio Map：
\begin{equation}
R(u,v)=
\frac{\tilde{S}(u,v)-\tilde{S}_{\min}}
{\tilde{S}_{\max}-\tilde{S}_{\min}+\epsilon},
\label{eq:4_27_radiomap_normalization}
\end{equation}
其中，\(\tilde{S}_{\min}\) 和 \(\tilde{S}_{\max}\) 分别表示平滑矩阵中的最小值和最大值，\(\epsilon\) 为防止分母为零而引入的极小常数。式~\eqref{eq:4_24_grid_assignment}--式~\eqref{eq:4_27_radiomap_normalization} 给出了从离散 RSSI 样本到连续 Radio Map 的基本构造过程。

\section{基于知识蒸馏的鲁棒多无人机协同定位方法}
\label{sec:4_4}

\subsection{多无人机掉点场景下的协同定位背景与问题引入}
\label{subsec:4_3_1}

在实际部署中，多无人机协同定位并不总能保持完整拓扑。受续航衰减、链路中断、局部遮挡和电磁干扰等因素影响，部分无人机可能在任务过程中暂时失效，从而使原本完整的协同观测退化为残缺观测。对于多目标定位系统而言，这意味着问题不再是固定节点规模下的静态估计，而是动态拓扑退化条件下的鲁棒定位。

从场景形式上看，本节所考虑的是：多架无人机在目标区域上空或周边以分布式方式执行协同扫描，不同节点从各自空间视角对地面多个终端进行 RSSI 观测，并将测量结果汇聚到统一定位模型中完成联合判别。在理想条件下，这些无人机节点能够共同形成较为完整的空间覆盖，使系统同时利用多视角观测、节点间几何分布以及区域响应差异来识别目标位置。然而，一旦部分节点在任务过程中退出协同网络，原本由多节点共同支撑的空间观测结构便会被破坏。

需要进一步指出的是，节点掉点并不是后续算法阶段才出现的附加问题，而是多无人机协同扫描场景本身的重要任务约束。受续航衰减、通信中断、局部遮挡以及电磁干扰等因素影响，部分无人机可能在任务执行过程中暂时失效或退出协同链路，从而使原本完整的多节点空间观测退化为不完整观测，并直接影响目标区域的空间覆盖结构以及后续 Radio Map 与区域级判别结果。

为对这一现象作抽象描述，设协同观测网络由 \(N_u\) 架无人机构成，其节点可用状态记为
\begin{equation}
\boldsymbol{\delta}
=
[\delta_1,\delta_2,\ldots,\delta_{N_u}]^{T},
\qquad
\delta_i\in\{0,1\},
\label{eq:4_2_uav_availability_vector}
\end{equation}
其中，\(\delta_i=1\) 表示第 \(i\) 架无人机处于有效观测状态，\(\delta_i=0\) 表示该节点在当前时刻不可用。于是，协同观测网络的有效覆盖率可定义为
\begin{equation}
\rho_{\mathrm{cov}}
=
\frac{1}{N_u}\sum_{i=1}^{N_u}\delta_i,
\label{eq:4_3_effective_coverage_ratio}
\end{equation}
其取值越小，说明有效观测节点越少，目标区域可获得的空间观测越不完整。当 \(\rho_{\mathrm{cov}}\) 随任务过程动态变化时，多目标定位问题便不再只是静态稀疏采样问题，而进一步演化为同时受到采样稀疏性与协同拓扑退化共同影响的联合判别问题。

这种破坏并不仅仅表现为少了几个测量值。对于多目标并发定位而言，节点掉点还会带来三个直接后果：其一，目标区域的空间采样密度下降，部分区域可能由原本的多视角覆盖退化为少量节点观测；其二，节点间的几何基线被削弱，不同目标之间原本可分离的空间响应更容易发生混叠；其三，后续 Radio Map 与区域判别所依赖的整体纹理结构会随有效节点数量变化而发生失真。因此，掉点问题本质上是观测信息量、空间覆盖关系和协同几何结构的同步退化。

设时刻 \(t\) 的有效节点集合记为
\begin{equation}
\mathcal{V}_{t}^{\mathrm{act}}
=
\{u_i \mid u_i \text{ 在时刻 } t \text{ 可用}\},
\qquad
N_t^{\mathrm{act}} = |\mathcal{V}_{t}^{\mathrm{act}}|,
\label{eq:4_40_active_node_set}
\end{equation}
则对应的拓扑退化比例可写为
\begin{equation}
\eta_t
=
1-\frac{N_t^{\mathrm{act}}}{N},
\label{eq:4_41_dropout_ratio}
\end{equation}
其中，\(N\) 表示理想状态下的总节点数。随着 \(\eta_t\) 增大，系统可获得的空间观测视角不断减少，输入特征的维度与结构也随之变化，这会直接削弱多目标定位中对空间分布形态的刻画能力。特别是在中高掉点比例条件下，即使目标真实分布不变，定位模型接收到的空间响应图样也会与完整网络条件下产生明显偏离。

对于传统依赖固定维度输入的定位模型而言，节点掉点通常只能通过补零、截断或丢弃样本等方式处理。这类做法虽然能够在形式上维持输入结构，但无法恢复缺失节点原本提供的空间感知信息，也难以有效利用残存节点之间仍然存在的协同约束关系，因此定位性能往往会随掉点程度上升而快速退化。基于这一认识，本章进一步考虑：如何在随机掉点条件下，仍尽可能保持接近完整网络的定位能力。

从方法设计角度看，仅依赖残缺观测本身进行监督学习，其性能上界天然受限于缺失后的信息量。为此，本文引入教师—学生知识蒸馏框架：教师模型在完整节点条件下学习全局空间映射关系，学生模型在随机掉点条件下完成定位推断，并通过蒸馏监督逼近完整网络的判别能力。由此，后续小节将围绕完整观测知识如何迁移到残缺观测条件下展开具体建模。

\subsection{知识蒸馏基本原理与掉点鲁棒定位建模}
\label{subsec:4_3_2}

针对上一小节所述的多无人机随机掉点问题，本文引入知识蒸馏（Knowledge Distillation）思想，对残缺观测条件下的协同定位模型进行鲁棒性增强。我们认为，虽然掉点后的无人机网络只能获得部分节点观测，但在离线训练阶段，系统仍然可以利用完整无人机网络所形成的全量观测数据，学习更充分的空间几何映射关系。若能够将这种由完整网络学得的判别能力迁移到掉点网络中，则有望在不增加额外硬件部署成本的前提下，缓解节点缺失所带来的性能退化。

从一般意义上看，知识蒸馏是一种教师—学生学习框架。教师模型通常具有更完整的输入视角或更强的表征能力，其输出不仅包含最终类别信息，还隐含了不同类别之间的相对相似性；学生模型则在输入受限或结构更轻量的条件下，通过模仿教师模型的输出分布来获得额外监督。与仅使用硬标签训练相比，蒸馏学习所利用的软标签能够保留类别间的细粒度关系，因此更适合在信息不完整条件下指导模型逼近高质量决策边界。

结合本章问题，本文考虑一个由 \(N\) 架无人机组成的协同感知网络。对于地面终端 \(j\)，第 \(i\) 架无人机的空间位置记为 \((x_i,y_i,z_i)\)，对应 RSSI 测量值记为 \(R_{i,j}\)。为了同时表征测量值及其空间来源，本文将单节点感知特征定义为
\begin{equation}
X_{i,j}=[x_i,y_i,z_i,R_{i,j}].
\label{eq:4_44_single_node_feature}
\end{equation}
\noindent 设待定位区域被离散为 \(M\) 个栅格类别，则定位任务可表示为栅格级多分类问题，模型输出为目标属于各候选栅格的概率分布。

在理想状态下，所有无人机节点均正常工作，此时教师网络接收完整观测并输出全局判别分布，可表示为
\begin{equation}
\boldsymbol{\pi}^{\,T}
=
f_T(\{X_{i,j}\}_{i=1}^{N}),
\label{eq:4_48_teacher_output}
\end{equation}
其中，\(f_T(\cdot)\) 表示教师网络，\(\boldsymbol{\pi}^{\,T}\in\mathbb{R}^{M}\) 为教师网络输出的栅格概率分布。

在实际任务中，由于部分无人机可能掉点，学生网络仅能利用残缺观测完成定位，其输出可写为
\begin{equation}
\boldsymbol{\pi}^{\,S}
=
f_S(\{X_{i,j}\mid S_i=1\}),
\label{eq:4_51_student_output}
\end{equation}
其中，\(S_i\in\{0,1\}\) 表示节点状态，\(S_i=1\) 表示第 \(i\) 架无人机工作正常，\(S_i=0\) 表示该节点已掉点；\(f_S(\cdot)\) 表示学生网络，\(\boldsymbol{\pi}^{\,S}\in\mathbb{R}^{M}\) 为学生网络对各栅格的预测概率分布。与教师网络相比，学生网络的输入由完整节点集合退化为残缺节点集合，其判别能力因此更容易受观测缺失影响。

在仅使用真实标签监督的情况下，学生网络通过交叉熵损失完成训练。设真实栅格标签对应的 one-hot 向量为 \(\mathbf{y}\)，则常规监督损失可写为
\begin{equation}
\mathcal{L}_{\mathrm{ce}}
=
-\sum_{m=1}^{M} y_m \log \pi^{S}_{m}.
\label{eq:4_52_ce}
\end{equation}
然而，当输入仅由残缺观测构成时，学生网络可利用的信息量显著低于教师网络，若仍仅依赖硬标签训练，则其学习过程容易受限于局部最优解，难以逼近完整网络条件下的判别性能。

为此，本文引入蒸馏监督。设 \(\ell\in\{T,S\}\) 表示教师网络或学生网络，\(z_m^{\ell}\) 表示其在第 \(m\) 个栅格上的 logits 输出，则对应的软标签分布统一写为
\begin{equation}
q^{\ell}_{m}
=
\frac{\exp(z^{\ell}_{m}/T_d)}
{\sum_{n=1}^{M}\exp(z^{\ell}_{n}/T_d)},
\qquad \ell\in\{T,S\},
\label{eq:4_53_soft_targets}
\end{equation}
其中，\(T_d\) 为温度系数。相比于 one-hot 标签，软标签分布不仅保留了最优类别信息，还反映了目标在不同栅格之间的相对相似性。对于本章所研究的地理栅格定位任务而言，这种信息尤为重要，因为相邻栅格之间往往具有空间连续性和几何相关性，而完整教师网络正是通过全量节点观测学习到了这种全局空间映射规律。

基于此，本文采用 Kullback--Leibler 散度度量教师与学生输出分布之间的差异，其蒸馏损失定义为
\begin{equation}
\mathcal{L}_{\mathrm{kd}}
=
T_d^2
\sum_{m=1}^{M}
q^{T}_{m}
\log
\frac{q^{T}_{m}}{q^{S}_{m}}.
\label{eq:4_55_kd}
\end{equation}
最终，学生网络的训练目标由硬标签监督与蒸馏监督共同组成，即
\begin{equation}
\mathcal{L}
=
\alpha \mathcal{L}_{\mathrm{ce}}
+
\beta \mathcal{L}_{\mathrm{kd}},
\label{eq:4_56_total_loss}
\end{equation}
其中，\(\alpha\) 和 \(\beta\) 分别表示两类损失项的权重系数。

式~\eqref{eq:4_56_total_loss} 所对应的训练目标可以作如下理解：\(\mathcal{L}_{\mathrm{ce}}\) 保证学生网络能够利用真实栅格标签完成基本分类定位任务，\(\mathcal{L}_{\mathrm{kd}}\) 则进一步迫使学生网络在面对残缺观测时，去模仿完整教师网络所形成的全局判别逻辑。这样一来，学生模型学习到的就不再只是残缺输入条件下如何拟合硬标签，而是在完整观测视角下，该目标更可能位于哪些栅格、不同栅格之间的概率关系如何变化。这正是知识蒸馏能够用于掉点鲁棒定位的根本原因。

综上所述，本文将多无人机掉点场景下的协同定位问题建模为一个基于知识迁移的多分类学习问题：教师网络利用完整无人机网络构建的 Radio Map 学习全局空间几何映射，学生网络则在随机掉点条件下通过模仿教师输出分布，提高残缺观测条件下的定位鲁棒性。基于上述建模，下一小节将进一步给出教师网络与学生网络的具体架构设计。

\subsection{网络架构设计}
\label{subsec:4_3_3}

为实现完整无人机网络知识向掉点网络的有效迁移，本文构建了一个面向多无人机协同定位任务的教师—学生蒸馏架构。该架构以图神经网络为基础，将多无人机空间观测组织为动态图输入，并分别在完整节点条件和残缺节点条件下完成全局特征提取与栅格分类。其核心目标在于：教师网络在全量节点观测下学习完整的空间几何映射规律，学生网络在随机掉点条件下利用残存节点信息完成定位推断，并通过蒸馏损失逼近教师网络的输出分布。

图~\ref{fig:kd_teacher_student_architecture} 给出了本文教师—学生网络架构示意图。图中左侧为完整无人机网络输入下的教师模型，右侧为随机掉点后的学生模型，中间为蒸馏损失与监督关系。该图对应展示了完整图输入、残缺图输入、动态建图、图注意力特征提取、全局池化读出以及分类输出等关键模块之间的连接关系。

\vspace{-0.2em}
\begin{figure}[H]
    \centering
    \includegraphics[width=0.98\textwidth]{ch4/KD网络架构图.png}
    \caption{教师—学生网络架构}
    \label{fig:kd_teacher_student_architecture}
\end{figure}
\vspace{-0.4em}

从输入形式上看，教师网络与学生网络对应两类不同的观测状态。教师网络接收完整无人机网络所形成的全量节点特征集合
\begin{equation}
\mathcal{X}_{\mathrm{complete}}
=
\{X_{1,j},X_{2,j},\dots,X_{N,j}\},
\end{equation}
其中，单节点特征 \(X_{i,j}=[x_i,y_i,z_i,R_{i,j}]\) 由无人机空间位置与对应 RSSI 测量值共同构成。学生网络接收随机掉点条件下的残缺节点集合
\begin{equation}
\mathcal{X}_{\mathrm{missing}}
=
\{X_{i,j}\mid S_i=1\},
\end{equation}
其中，\(S_i\in\{0,1\}\) 表示节点状态，掉点后有效节点数由 \(N\) 退化为 \(N-k\)。两类输入在观测完整性上存在显著差异，这也是后续蒸馏学习的基础。

考虑到多无人机协同定位问题中各节点之间天然存在空间拓扑关系，本文不采用简单的向量拼接方式组织输入，而是将观测数据建模为图结构。具体而言，每一架无人机对应图中的一个节点，节点特征由其三维坐标与 RSSI 测量共同给出。对于教师网络，图中包含全部 \(N\) 个节点；对于学生网络，图中仅保留当前未掉点的 \(N-k\) 个有效节点。为充分利用节点间的协同关系，本文在图构建阶段采用全连接策略，使任意两个有效节点之间都建立边连接。这样处理的目的在于避免人为限制节点间的信息传播范围，使网络能够在训练过程中自主学习不同空间位置之间的相关性强弱。若记输入图为 \(\mathcal{T}=(\mathcal{V},\mathcal{E})\)，节点 \(i\) 的输入特征为 \(\mathbf{h}_i\)，则对于任意相邻节点 \(j\in\mathcal{N}(i)\)，图注意力层首先计算节点对之间的未归一化注意力分数
\begin{equation}
e_{ij}
=
\mathrm{LeakyReLU}
\!\left(
\mathbf{a}^{T}
\left[
\mathbf{W}\mathbf{h}_{i}
\;\Vert\;
\mathbf{W}\mathbf{h}_{j}
\right]
\right),
\label{eq:4_57_gat_score}
\end{equation}
其中，\(\mathbf{W}\) 表示共享线性变换矩阵，\(\mathbf{a}\) 表示注意力权重向量，\(\Vert\) 表示向量拼接操作。随后，对节点 \(i\) 的所有邻居进行 softmax 归一化，得到注意力系数
\begin{equation}
\alpha_{ij}
=
\frac{\exp(e_{ij})}
{\sum_{k\in\mathcal{N}(i)}\exp(e_{ik})},
\label{eq:4_58_gat_softmax}
\end{equation}
并据此完成节点特征更新：
\begin{equation}
\mathbf{h}_{i}^{\prime}
=
\sigma\!\left(
\sum_{j\in\mathcal{N}(i)}
\alpha_{ij}\mathbf{W}\mathbf{h}_{j}
\right),
\label{eq:4_59_gat_update}
\end{equation}
其中，\(\sigma(\cdot)\) 表示非线性激活函数。式~\eqref{eq:4_57_gat_score}--式~\eqref{eq:4_59_gat_update} 给出了本文所采用图注意力机制的基本计算形式。

在图结构建立后，教师网络与学生网络均采用多层图注意力网络作为核心特征提取模块。图注意力层的作用在于，根据节点之间的特征相关性自适应分配邻域权重，从而在特征传播过程中突出更有判别价值的节点贡献。对于教师网络而言，多层图注意力结构能够在全量节点条件下学习完整的协同感知关系，并形成对全局空间布局较为充分的表征；对于学生网络而言，同样的结构则能够在残缺拓扑条件下挖掘剩余节点之间仍然存在的空间补偿关系，从而缓解局部观测缺失带来的信息损失。

设经过图注意力层堆叠后，第 \(i\) 个节点的高层表示记为 \(\mathbf{h}_i\)，则教师网络在图级读出阶段采用全局平均池化，将全部节点特征汇聚为固定维度的图表示
\begin{equation}
\mathbf{h}_{G}^{T}
=
\frac{1}{N}\sum_{i=1}^{N}\mathbf{h}_i.
\end{equation}
对应地，学生网络在残缺节点条件下得到的图级表示为
\begin{equation}
\mathbf{h}_{G}^{S}
=
\frac{1}{N-k}\sum_{i\in\mathcal{V}_{\mathrm{act}}}\mathbf{h}_i,
\end{equation}
其中，\(\mathcal{V}_{\mathrm{act}}\) 表示当前有效节点集合。采用全局平均池化的原因在于，该读出方式能够在保持输出维度固定的同时，自然适应输入节点数量的变化，为后续分类头提供统一形式的图级特征。

在读出层之后，教师网络与学生网络均采用多层感知机（Multi-Layer Perceptron, MLP）完成最终的栅格分类。教师网络的输出分类得分（logits）记为
\begin{equation}
\mathbf{z}^{T}
=
\mathrm{MLP}(\mathbf{h}_{G}^{T}),
\end{equation}
学生网络的输出 logits 记为
\begin{equation}
\mathbf{z}^{S}
=
\mathrm{MLP}(\mathbf{h}_{G}^{S}).
\end{equation}
两者的输出维度均为 \(M\)，对应目标区域离散后的 \(M\) 个地理栅格类别。教师网络在完整输入条件下生成全局判别分布，学生网络则在残缺输入条件下给出对应的栅格概率预测。

在训练阶段，教师网络的 logits 经过温度系数 \(T_d\) 的 softmax 变换后，生成软标签分布，用于向学生网络提供额外监督；学生网络则同时保留两类输出，一类用于与真实栅格标签计算交叉熵损失，另一类用于与教师网络输出进行蒸馏匹配。由此，整个架构形成了完整观测指导残缺观测、教师网络监督学生网络的训练闭环。需要强调的是，这里教师模型与学生模型在总体结构上保持一致，但二者的输入视角不同：前者对应完整图输入，后者对应掉点后的残缺图输入。正是这种输入不对称性，使蒸馏学习能够在不增加额外硬件成本的情况下，将完整网络中学得的全局空间判别能力迁移到掉点网络中。

从结构设计角度看，本文所构建的教师—学生网络具有三点特点。其一，采用图结构建模多无人机观测，避免了固定维度输入对节点完整性的严格依赖；其二，利用图注意力机制增强节点间空间相关性的自适应建模能力，使网络在掉点条件下仍可充分利用残存节点信息；其三，通过图级读出和统一分类头，将不同规模的节点输入映射为一致的栅格分类输出，从而为后续蒸馏训练提供了结构基础。

综上，本文教师—学生网络架构并不是两个独立分类器的简单并列，而是一个围绕完整观测知识向残缺观测迁移而构建的协同学习框架。教师网络负责在全量节点条件下学习完整的空间几何映射规律，学生网络则在随机掉点条件下，通过图注意力传播、图级特征读出和蒸馏监督，尽可能恢复完整网络条件下的定位能力。基于该架构，下一小节将进一步给出面向残缺观测的蒸馏训练策略。

\section{实验结果分析}
\label{sec:4_5}

\subsection{实验设置}
\label{subsec:4_4_1}

为验证所提基于知识蒸馏的掉点鲁棒多无人机协同定位方法，本节在复杂城市环境下构建多无人机协同定位实验平台，并在完整观测、随机掉点和蒸馏增强三种条件下开展对比实验。整个实验流程包括场景建模、数据集生成、残缺观测构造、教师网络训练、学生网络蒸馏训练以及测试分析等环节。

实验场景采用基于真实地理信息数据构建的复杂城市街区。目标区域大小设定为 \(500\,\mathrm{m}\times 500\,\mathrm{m}\)，并均匀划分为 \(20\times 20\) 个离散栅格，对应 \(M=400\) 个空间类别。经此离散化处理后，地面终端定位问题可表述为栅格级多分类任务，模型输出为目标位于各候选栅格的概率分布。

协同感知网络由 \(N=16\) 架无人机构成。每架无人机在给定时刻均具有已知三维空间坐标，并可对地面终端发射信号进行 RSSI 测量。对于任一目标终端，单节点观测特征定义为
\begin{equation}
X_{i,j}=[x_i,y_i,z_i,R_{i,j}].
\end{equation}
完整网络条件下，系统可获得 16 架无人机联合形成的全量观测；当部分节点掉点时，仅保留有效节点观测，构成残缺输入样本。

数据集由复杂城市传播场建模结果生成，并进一步完成 RSSI 栅格映射与 Radio Map 生成。完整观测条件下形成的样本用于教师网络训练，随机掉点后的残缺观测用于学生网络训练。为评估模型在拓扑退化条件下的稳定性，实验设置五种掉点比例，分别为 \(0\%\)、\(12.5\%\)、\(25\%\)、\(37.5\%\) 和 \(50\%\)，对应掉点节点数分别为 0、2、4、6 和 8。训练阶段与测试阶段均通过随机采样不同掉点组合构造残缺观测样本，以减小偶然拓扑分布对结果的影响。

对比方法共包括四类。其一为完整观测条件下的教师网络（Teacher），用作性能上界参考；其二为不采用知识蒸馏（Knowledge Distillation, KD）的学生网络（Student-NoKD）；其三为传统插值补全后再执行定位的基线方法（Interp-Loc）；其四为本文提出的蒸馏增强学生网络（Student-KD）。上述设置分别对应完整观测上界、无蒸馏残缺学习、传统补全思路以及蒸馏增强方案，能够较完整地体现知识蒸馏在掉点场景中的作用。

教师模型与学生模型均采用三层图注意力网络（Graph Attention Network, GAT）作为核心特征提取模块。每层隐藏维度设为 64，注意力头数设为 4，激活函数采用修正线性单元（Rectified Linear Unit, ReLU），Dropout 系数设为 0.2。图级读出阶段采用全局平均池化，将变长节点集合映射为固定维度的图表示；分类头由两层多层感知机构成，结构为 \(64 \rightarrow 128 \rightarrow 400\)，最终输出 400 维栅格分类 logits。

网络训练采用 Adam 优化器。初始学习率设为 \(1\times10^{-3}\)，权重衰减系数设为 \(1\times10^{-4}\)，批大小设为 64，最大训练轮数设为 800。学习率衰减采用 StepLR 分段衰减策略，每 200 个 epoch 衰减为原来的 0.5。同时引入早停机制（early stopping），当验证集性能在连续 50 个 epoch 内无明显提升时终止训练。知识蒸馏部分的温度系数设为 \(T_d=4\)，交叉熵损失与蒸馏损失的权重分别设为 \(\alpha=1.0\) 和 \(\beta=0.7\)。

全部样本按照 70\%、15\%、15\% 的比例划分为训练集、验证集和测试集。其中，训练集用于教师网络与学生网络的参数优化，验证集用于模型选择与早停判断，测试集用于最终性能评估。

性能评价采用平均定位误差（Mean Error）、均方根误差（RMSE）以及 90\% 分位误差（90\% Error）三项指标。设真实终端位置为 \(\mathbf{p}^{\mathrm{gt}}=[x^{\mathrm{gt}},y^{\mathrm{gt}}]^T\)，预测栅格中心位置为 \(\hat{\mathbf{p}}=[\hat{x},\hat{y}]^T\)，则单样本定位误差定义为 \(e=\left\|\hat{\mathbf{p}}-\mathbf{p}^{\mathrm{gt}}\right\|_2\)。其中，Mean Error 反映整体平均定位偏差，RMSE 用于衡量总体误差水平，90\% Error 用于刻画误差分布尾部特性。针对掉点鲁棒性分析，本文进一步考察不同掉点比例下 Mean Error 的变化趋势。

表~\ref{tab:ch4_hyperparameters} 给出了本文实验所采用的主要参数设置。

\begingroup
\setlength{\LTleft}{\fill}
\setlength{\LTright}{\fill}
\begin{longtable}{cc}
    \caption{实验参数设置表}
    \label{tab:ch4_hyperparameters}\\
    \toprule
    \textbf{参数名称} & \textbf{参数取值} \\
    \midrule
    \endfirsthead
    \caption[]{实验参数设置表（续）}\\
    \toprule
    \textbf{参数名称} & \textbf{参数取值} \\
    \midrule
    \endhead
    目标区域大小 & \(500\,\mathrm{m}\times500\,\mathrm{m}\) \\
    栅格划分 & \(20\times20\) \\
    栅格类别数 \(M\) & 400 \\
    无人机数量 \(N\) & 16 \\
    掉点比例 & \(0\%,12.5\%,25\%,37.5\%,50\%\) \\
    GAT 层数 & 3 \\
    隐藏维度 & 64 \\
    注意力头数 & 4 \\
    激活函数 & ReLU \\
    Dropout & 0.2 \\
    图级读出方式 & Global Average Pooling \\
    分类头结构 & \(64 \rightarrow 128 \rightarrow 400\) \\
    最大训练轮数 & 800 \\
    批大小 & 64 \\
    初始学习率 & \(1\times10^{-3}\) \\
    优化器 & Adam \\
    权重衰减 & \(1\times10^{-4}\) \\
    学习率衰减策略 & StepLR \\
    衰减步长 & 200 epoch \\
    衰减系数 & 0.5 \\
    Early stopping patience & 50 \\
    温度系数 \(T_d\) & 4 \\
    损失权重 \(\alpha\) & 1.0 \\
    损失权重 \(\beta\) & 0.7 \\
    数据集划分 & 70\% / 15\% / 15\% \\
    \bottomrule
\end{longtable}
\endgroup

为补充说明不同方法在工程实现层面的计算代价，本文进一步从单样本推理复杂度与额外训练开销两个角度，对第四章涉及的主要对比方法进行分析。设完整协同网络节点数为 \(N\)，掉点后有效节点数为 \(\tilde{N}=N-k\)，图注意力层数为 \(L_g\)，注意力头数为 \(H\)，隐藏维度为 \(d\)，输出栅格类别数为 \(M\)，插值重构的网格规模为 \(G\)。由于本文采用全连接图建模，多头图注意力网络单层计算复杂度可近似记为 \(O(HN^2d)\) 或 \(O(H\tilde{N}^2d)\)。对于传统插值补全方法，记其下游定位器的推理复杂度为 \(C_{\mathrm{loc}}\)。据此，不同方法的复杂度比较结果如表~\ref{tab:ch4_complexity_compare} 所示。

\begin{table}[H]
    \centering
    \caption{第四章不同方法的复杂度比较}
    \label{tab:ch4_complexity_compare}
    \begin{tabular}{lccc}
        \toprule
        \textbf{方法} & \textbf{主要计算模块} & \textbf{单样本推理复杂度} & \textbf{额外训练开销} \\
        \midrule
        Teacher & 完整图 GAT + MLP & \(O(L_gHN^{2}d+dM)\) & 无 \\
        Student-NoKD & 残缺图 GAT + MLP & \(O(L_gH\tilde{N}^{2}d+dM)\) & 无 \\
        Interp-Loc & 插值重构 + 下游定位器 & \(O(G)+C_{\mathrm{loc}}\) & 无 \\
        Student-KD & 残缺图 GAT + MLP & \(O(L_gH\tilde{N}^{2}d+dM)\) & \(O(L_gHN^{2}d+M)\) \\
        \bottomrule
    \end{tabular}
\end{table}

由表~\ref{tab:ch4_complexity_compare} 可见，Student-KD 的额外开销主要体现在训练阶段需要引入教师网络前向计算与蒸馏监督，而其测试阶段的推理复杂度与 Student-NoKD 保持一致。相比之下，Interp-Loc 需要额外执行插值重构步骤，因此其总体开销还受到重构网格规模 \(G\) 的影响。在本文实验设置下，完整协同网络节点数为 \(N=16\)，当掉点比例分别为 \(12.5\%\)、\(25\%\)、\(37.5\%\) 和 \(50\%\) 时，Student 类方法对应的有效节点数 \(\tilde{N}\) 分别为 14、12、10 和 8，其图传播主项相对完整网络约降为 \(76.6\%\)、\(56.3\%\)、\(39.1\%\) 和 \(25.0\%\)。

\subsection{性能对比实验}
\label{subsec:4_4_2}

为比较不同方法在残缺观测条件下的定位能力，本文首先开展性能对比实验。在该实验中，固定测试阶段的掉点比例为 \(25\%\)，即随机屏蔽 4 架无人机的观测数据，以模拟中等程度的节点失效场景。在相同测试集和相同掉点模式下，分别对 Teacher、Student-NoKD、Interp-Loc 以及 Student-KD 四种方法进行测试，并统计各测试样本的定位误差分布。其中，Teacher 表示在完整观测条件下训练和测试的教师网络，用作完整节点输入下的性能上界；Student-NoKD 表示仅利用掉点后的残缺观测、且不采用知识蒸馏训练的学生网络；Interp-Loc 表示先对缺失节点观测进行插值补全，再执行后续定位判别的传统基线方法；Student-KD 则表示本文提出的蒸馏增强学生网络，即在残缺观测条件下利用教师软标签监督进行训练的鲁棒定位模型。

图~\ref{fig:ch4_error_cdf} 给出了四种方法在 \(25\%\) 掉点条件下的定位误差累计分布曲线。可以看出，Teacher 曲线整体位于最左侧，说明完整观测条件下的教师网络具有最优性能；Student-NoKD 曲线最靠右，表明仅依赖残缺观测直接训练的学生模型在误差分布上出现明显劣化。Interp-Loc 相较 Student-NoKD 有一定改善，说明传统插值补全策略能够在一定程度上缓解观测缺失带来的影响，但由于该方法并未显式利用完整网络所学习到的全局空间判别知识，因此性能提升仍然有限。相比之下，Student-KD 曲线整体更接近 Teacher，并显著位于 Student-NoKD 左侧，说明蒸馏监督能够有效缩小残缺网络与完整网络之间的性能差距。

\begin{figure}[htbp]
    \centering
    \includegraphics[width=0.82\textwidth]{ch4/9_不同模型平均定位误差对比图.png}
    \caption{不同方法在 \(25\%\) 掉点条件下的定位误差累计分布曲线}
    \label{fig:ch4_error_cdf}
\end{figure}

为进一步进行定量比较，表~\ref{tab:ch4_overall_compare} 给出了不同方法在 \(25\%\) 掉点条件下的误差统计结果，包括 Mean Error、RMSE 和 90\% Error 三项指标。该表从均值、整体波动和尾部误差三个维度对模型性能进行了量化评估。

\begin{table}[htbp]
    \centering
    \caption{不同方法在 \(25\%\) 掉点条件下的定位性能比较}
    \label{tab:ch4_overall_compare}
    \begin{tabular}{lccc}
        \toprule
        \textbf{方法} & \textbf{Mean Error (m)} & \textbf{RMSE (m)} & \textbf{90\% Error (m)} \\
        \midrule
        Teacher & 9.8 & 12.6 & 21.8 \\
        Student-NoKD & 22.4 & 27.9 & 46.5 \\
        Interp-Loc & 17.6 & 22.1 & 36.8 \\
        Student-KD & 12.8 & 16.0 & 26.4 \\
        \bottomrule
    \end{tabular}
\end{table}

由表~\ref{tab:ch4_overall_compare} 可见，Teacher 在完整观测条件下取得最低的 Mean Error、RMSE 和 90\% Error，体现了全量节点输入所带来的性能上界。Student-NoKD 的三项误差指标均显著增大，说明在掉点条件下仅依赖硬标签监督训练的学生模型难以恢复完整网络中的空间判别能力。Interp-Loc 在三项指标上均优于 Student-NoKD，表明对缺失观测进行插值补全确实具有一定作用，但其整体性能仍明显落后于 Student-KD。相比之下，Student-KD 的 Mean Error、RMSE 和 90\% Error 分别降至 12.8 m、16.0 m 和 26.4 m，较 Student-NoKD 均取得了约 40\% 的误差下降，并显著逼近 Teacher，说明知识蒸馏能够有效迁移完整网络中学习到的全局空间结构知识。

图~\ref{fig:ch4_error_bar} 给出了不同方法在 \(25\%\) 掉点条件下的 Mean Error 与 90\% Error 统计结果。结合图~\ref{fig:ch4_error_cdf} 与表~\ref{tab:ch4_overall_compare} 可以看出，Student-KD 不仅在平均误差上明显优于 Student-NoKD 和 Interp-Loc，而且在尾部大误差样本上也表现出更强的抑制能力。这说明蒸馏方法的作用并非仅体现在整体均值改善上，同时也提高了模型在复杂样本条件下的稳定性。

\begin{figure}[htbp]
    \centering
    \includegraphics[width=0.82\textwidth]{ch4/10_不同模型九十分位误差对比图.jpg}
    \caption{不同方法在 \(25\%\) 掉点条件下的误差统计柱状图}
    \label{fig:ch4_error_bar}
\end{figure}

综合该实验可以得到两个结论：其一，残缺观测会显著削弱学生模型的定位能力，Student-NoKD 与 Teacher 之间存在明显性能差距；其二，相较于仅依赖传统插值补全的 Interp-Loc，Student-KD 能够更充分地利用教师网络输出的软标签信息，从而在平均误差和尾部误差两个层面均取得更优结果。这表明，知识蒸馏确实能够将完整网络中学习到的全局空间结构知识迁移到掉点条件下的学生模型中。

\subsection{掉点鲁棒性实验}
\label{subsec:4_4_3}

为进一步验证所提方法在不同拓扑退化条件下的稳定性，本文开展掉点鲁棒性实验。在该实验中，分别设置 \(0\%\)、\(12.5\%\)、\(25\%\)、\(37.5\%\) 和 \(50\%\) 五种掉点比例，并在每种掉点比例下测试 Student-NoKD 与 Student-KD 的定位精度。实验重点关注掉点比例增加时两类学生网络性能退化的快慢差异。

随着掉点比例逐步增加，两类学生网络的定位误差均呈上升趋势，这与有效观测节点数量减少所导致的信息损失相一致。然而，两者的退化速率存在明显差异。Student-NoKD 完全依赖残缺观测进行判别，其误差会随掉点比例上升而快速增大；而 Student-KD 在教师网络软标签监督下保留了更多完整网络中的全局空间知识，因此误差增长更为平缓。特别是在中高掉点比例条件下，两类方法之间的性能差距进一步扩大，说明蒸馏方法在严重拓扑退化条件下能够更充分地体现优势。

为进一步量化鲁棒性差异，表~\ref{tab:ch4_dropout_compare} 给出了不同掉点比例条件下两类学生网络的 Mean Error 比较结果。

\begin{table}[htbp]
    \centering
    \caption{不同掉点比例下 Student-NoKD 与 Student-KD 的误差比较}
    \label{tab:ch4_dropout_compare}
    \begin{tabular}{lcc}
        \toprule
        \textbf{掉点比例} & \textbf{Student-NoKD Mean Error (m)} & \textbf{Student-KD Mean Error (m)} \\
        \midrule
        \(0\%\) & 13.6 & 11.2 \\
        \(12.5\%\) & 17.8 & 11.9 \\
        \(25\%\) & 22.4 & 12.8 \\
        \(37.5\%\) & 29.7 & 14.6 \\
        \(50\%\) & 38.9 & 17.3 \\
        \bottomrule
    \end{tabular}
\end{table}

由表~\ref{tab:ch4_dropout_compare} 可以看出，在无掉点条件下，Student-KD 与 Student-NoKD 的性能差距相对有限；但随着掉点比例不断升高，二者之间的差距逐步扩大。当掉点比例达到 \(50\%\) 时，Student-NoKD 的 Mean Error 已上升至 38.9 m，而 Student-KD 仍保持在 17.3 m，说明蒸馏增强学生网络在高掉点条件下仍具有较强的误差控制能力。总体来看，本文方法并非仅在单一掉点设置下有效，而是在较宽范围的网络退化条件下均表现出更好的鲁棒性。这说明知识蒸馏所提供的并不是某一固定模式下的经验补偿，而是更具有普适性的全局空间先验迁移能力。

% \subsection{定位结果可视化分析}
% \label{subsec:4_4_4}

% 为更直观地展示教师网络与学生网络在空间定位上的差异，本文进一步选取若干典型测试样本进行可视化分析。对于每个样本，分别给出 Ground Truth、Teacher、Student-NoKD 和 Student-KD 的空间预测结果，以比较不同模型在同一残缺观测条件下的定位表现。

% 图~\ref{fig:ch4_grid_prediction_visual} 给出了典型测试样本的栅格预测可视化结果。图中每一行对应一个测试样本，每一列依次展示真实位置、教师网络预测、未蒸馏学生预测以及蒸馏后学生预测。通过这种方式，可以从单样本层面比较不同模型对目标空间分布的恢复能力。

% \begin{figure}[htbp]
%     \centering
%     \vspace{6.0cm}
%     \caption{典型测试样本的空间预测结果可视化对比}
%     \label{fig:ch4_grid_prediction_visual}
% \end{figure}

% 从可视化结果看，Teacher 在完整观测条件下通常能够较准确地恢复目标所在栅格，并给出相对集中的概率响应区域。Student-NoKD 在掉点条件下更容易出现响应区域偏移、概率分布扩散以及相邻栅格混淆等现象，尤其当目标位于边界区域或残存节点观测几何较差时，这种偏差更为明显。相比之下，Student-KD 的预测结果通常更接近 Teacher 与 Ground Truth，表现为响应区域更集中、峰值位置偏移更小。这说明蒸馏学习不仅改善了整体误差指标，也提升了学生网络在单样本层面的空间恢复质量。

% 进一步结合前文误差统计结果可以发现，Student-KD 相对于 Student-NoKD 的优势并非仅体现在均值意义上的误差下降，更体现在其对空间响应结构的保持能力上。也就是说，知识蒸馏不仅改善了模型对最终类别标签的判别结果，还增强了学生网络在残缺观测条件下对完整空间分布形态的恢复能力。

% \subsection{蒸馏机制消融实验}
% \label{subsec:4_4_5}

% 为进一步验证知识蒸馏机制在本文方法中的独立作用，本文开展消融实验。具体而言，分别比较仅采用交叉熵监督训练的学生模型与同时采用交叉熵监督和蒸馏监督训练的学生模型，并进一步考察温度系数与损失权重变化对结果的影响。消融实验统一在 \(25\%\) 掉点条件下进行，以保证与前述主实验设置一致。

% 表~\ref{tab:ch4_ablation_results} 给出了不同训练策略下学生模型的性能差异。

% \begin{table}[htbp]
%     \centering
%     \caption{蒸馏机制消融实验结果}
%     \label{tab:ch4_ablation_results}
%     \begin{tabular}{lccc}
%         \toprule
%         \textbf{模型设置} & \textbf{Mean Error (m)} & \textbf{RMSE (m)} & \textbf{90\% Error (m)} \\
%         \midrule
%         仅交叉熵监督 & 22.4 & 27.9 & 46.5 \\
%         交叉熵 + 蒸馏监督 & 12.8 & 16.0 & 26.4 \\
%         温度系数 \(T_d=1\) & 15.6 & 19.5 & 31.0 \\
%         损失权重 \(\beta=0.3\) & 14.8 & 18.2 & 29.4 \\
%         \bottomrule
%     \end{tabular}
% \end{table}

% 从表~\ref{tab:ch4_ablation_results} 可以看出，完整蒸馏方案在 Mean Error、RMSE 和 90\% Error 三项指标上均取得最优结果。仅采用交叉熵监督训练时，学生网络的 Mean Error、RMSE 和 90\% Error 分别达到 22.4 m、27.9 m 和 46.5 m，说明在掉点条件下仅依赖硬标签监督难以有效恢复完整网络中蕴含的全局空间判别知识。引入蒸馏监督后，三项指标分别下降至 12.8 m、16.0 m 和 26.4 m，表明知识蒸馏能够显著改善学生模型在残缺观测条件下的定位性能。

% 进一步地，当温度系数设置为 \(T_d=1\) 时，模型性能较完整蒸馏方案明显下降，说明适当的温度平滑对于保留教师输出中的类别相似性信息具有重要作用。当蒸馏损失权重减小为 \(\beta=0.3\) 时，模型性能同样出现一定退化，表明教师软标签监督在学生网络训练中不能过弱。总体来看，知识蒸馏并非附加性的训练技巧，而是本文所提出掉点鲁棒协同定位方法取得性能增益的关键来源之一。

\subsection{本节小结}
\label{subsec:4_4_6}

本节围绕性能对比、掉点鲁棒性与复杂度分析三个方面，对本文所提出的基于知识蒸馏的掉点鲁棒多无人机协同定位方法进行了实验验证。结果表明，在中等掉点条件下，Student-KD 的定位误差分布优于 Student-NoKD，并整体接近完整观测条件下的 Teacher；在不同掉点比例下，Student-KD 的性能退化曲线更平缓，说明其在随机掉点条件下具有更稳定的表现；复杂度分析则表明，蒸馏增强主要增加训练阶段开销，而测试阶段的推理复杂度与 Student-NoKD 保持一致。上述结果说明，本文方法能够在不增加额外硬件部署成本的前提下，减缓多无人机随机掉点所带来的定位性能退化问题。

\section{本章小结}
\label{sec:4_6}

本章面向复杂城市环境下的多无人机协同定位任务，围绕随机掉点条件下的鲁棒定位问题，提出了一种基于信号成像与知识蒸馏的协同定位方法。与传统依赖固定维度输入和完整节点观测的定位方法不同，本文重点考虑了多无人机网络在实际部署过程中可能出现的节点失效、链路中断和局部观测缺失现象，并在此基础上构建了面向残缺观测条件的鲁棒定位框架。

首先，针对复杂城市环境中射频传播机理复杂、离散观测难以直接用于定位的问题，本文基于真实地理信息数据和 MATLAB 工具链完成了复杂城市场景建模，并通过射线追踪获得了具有物理一致性的空间传播响应。在此基础上，进一步采用 RSSI 栅格映射与高斯平滑连续化表达方法，将无人机空间扫描获得的离散观测转化为具有区域结构的 Radio Map，为后续定位模型提供了统一输入表征。

其次，针对多无人机协同网络在随机掉点条件下输入维度不规则、传统刚性网络难以稳定处理的问题，本文将完整无人机网络生成的 Radio Map 作为教师输入，将掉点后的残缺 Radio Map 作为学生输入，构建了面向掉点鲁棒定位的教师—学生蒸馏框架。在该框架中，教师模型利用全量节点观测学习完整的空间几何映射关系，学生模型则在真实标签监督与软标签蒸馏监督的共同作用下，尽可能逼近完整网络条件下的定位能力。通过这种方式，本文在不增加额外硬件成本的前提下，实现了对缺失观测信息的有效补偿。

最后，实验结果表明，完整无人机网络能够为定位模型提供最充分的观测视角，而随机掉点会显著削弱未蒸馏学生网络的定位性能。相比之下，引入知识蒸馏后，学生模型在不同掉点比例条件下均表现出更稳定的分类精度和更小的定位误差，说明教师模型中学习到的全局空间判别规律能够有效迁移到残缺观测条件下的学生模型中。相关结果表明，本文方法能够在复杂场景中提升多无人机协同定位系统对节点失效的适应能力，增强系统的工程鲁棒性。

综上，本章从空间信号成像、残缺观测建模和知识迁移学习三个层面出发，完成了多无人机掉点条件下协同鲁棒定位方法的设计与验证。该方法为复杂环境中面向实际部署的多无人机协同定位提供了新的解决思路，也为后续进一步研究更大规模节点协同、更复杂掉点模式奠定了基础。
